% Vietnamese translation for OI Public Library
% Source-PDF: 比赛 CONTEST/2021“MINIEYE杯”中国大学生算法设计超级联赛/2021“MINIEYE杯”中国大学生算法设计超级联赛（5）/2021“MINIEYE杯”中国大学生算法设计超级联赛（5）-题解集.pdf
% Source-File: 比赛 CONTEST/2021“MINIEYE杯”中国大学生算法设计超级联赛/2021“MINIEYE杯”中国大学生算法设计超级联赛（5）/2021“MINIEYE杯”中国大学生算法设计超级联赛（5）-题解集.md
\documentclass[11pt]{article}
\usepackage{oipl-vn}

\title{Lời giải trận thứ năm HDU Multi-University 2021}
\author{}
\date{}

\begin{document}
\maketitle

\origtitle{HDU Multi-University 2021 MINIEYE Cup, Contest 5 Solutions}
\sourcepath{contest/hdu-2021-minieye-05-solutions.pdf}

\section{1001 Miserable Faith}

Với người quen thuộc link-cut tree, có thể nhận ra thao tác $1$ đang mô phỏng
thao tác \texttt{access} của LCT.

Vì vậy, khi thao tác $1$ làm thay đổi cạnh thật và cạnh ảo trong LCT, ta dùng
một cây đoạn để cập nhật theo đoạn, đồng thời duy trì đáp án của thao tác
$4$. Thao tác $2$ tương đương với truy vấn tại một điểm. Thao tác $3$ tương
đương với truy vấn tổng trên một cây con rồi trừ đi phần đáp án tại điểm
$u$.

Độ phức tạp là $O(n\log^2 n)$.

\section{1002 String Mod}

Ta có thể liệt kê số lần xuất hiện $x,y$ của hai ký tự $a,b$ trong chuỗi:
\[
\begin{aligned}
  \operatorname{ans}[i][j]
  &=
  \sum_{x=0}^{L}\sum_{y=0}^{L-x}
  [n\mid x-i][n\mid y-j]
  \binom{L}{x}\binom{L-x}{y}(k-2)^{L-x-y}.
\end{aligned}
\]
Dùng lọc nghiệm bằng căn đơn vị, với $w_n$ là căn bậc $n$ nguyên thủy:
\[
\begin{aligned}
  \operatorname{ans}[i][j]
  &=
  \frac1{n^2}
  \sum_{p=0}^{n-1}\sum_{q=0}^{n-1}
  \sum_{x=0}^{L}\sum_{y=0}^{L-x}
  w_n^{px}w_n^{qy}
  \binom{L}{x}\binom{L-x}{y}
  (k-2)^{L-x-y}
  w_n^{-pi}w_n^{-qj}  \\
  &=
  \frac1{n^2}
  \sum_{p=0}^{n-1}\sum_{q=0}^{n-1}
  (w_n^p+w_n^q+k-2)^Lw_n^{-pi}w_n^{-qj}.
\end{aligned}
\]
Đặt
\[
  A[i][p]=w_n^{-ip},\qquad
  B[p][q]=\frac1{n^2}(w_n^p+w_n^q+k-2)^L,\qquad
  C[q][j]=w_n^{-qj}.
\]
Khi đó
\[
  \operatorname{ans}=A\times B\times C.
\]
Độ phức tạp thời gian là $O(n^3+n^2\log L)$.

\section{1003 VC Is All You Need}

\paragraph{Phát biểu tương đương.}
Trong không gian $n$ chiều, cần tìm giá trị lớn nhất của $m$ sao cho có thể
tự chọn $m$ điểm, và với mọi cách tô hai màu các điểm đó, tức mọi trong
$2^m$ phương án tô màu, luôn tồn tại một siêu phẳng $n-1$ chiều tách nghiêm
ngặt hai màu.

\paragraph{Kết luận.}
\[
  m_{\max}=n+1.
\]

Do tính đơn điệu, chứng minh chia làm hai phần: chứng minh $m=n+1$ khả thi,
rồi chứng minh $m\ge n+2$ vô nghiệm.

Với phần thứ nhất, lấy
\[
  x_0=(0,0,\ldots,0),
\]
còn $x_i$ là véc-tơ $01$ có tọa độ thứ $i$ bằng $1$ và các tọa độ khác bằng
$0$, với $i\in[1,n]$. Xét hàm tuyến tính
\[
  f_w(x)=w_0+w_1x_1+\cdots+w_nx_n=w^T[1\ x]^T.
\]
Siêu phẳng phân loại là $f_w(x)\ge 0$.

Với một cách tô màu
\[
  S=\{(x_i,y_i)\},\qquad y_i\in\{+1,-1\},
\]
ta trực tiếp dựng $w$ sao cho $f_w(x_i)=y_i$ với mọi $i$. Hệ phương trình là
\[
  w_0=y_0,\qquad w_0+w_i=y_i\quad (i\in[1,n]),
\]
nên luôn giải được.

Với phần thứ hai, thay mỗi $x_i$ bởi $[1\ x_i]$. Trong không gian
$\mathbb{R}^{n+1}$, bất kỳ $n+2$ điểm nào cũng phụ thuộc tuyến tính. Do đó
tồn tại một tổ hợp tuyến tính không tầm thường bằng $0$. Chọn một hệ số khác
$0$ tương ứng với $x_0$, ta có thể viết
\[
  x_0=\sum_{i=1}^{n+1} a_ix_i.
\]
Ta dựng nhãn $y_i$ như sau: nếu $a_i\ge 0$ thì $y_i=1$, nếu $a_i<0$ thì
$y_i=-1$, và đặt $y_0=-1$.

Nếu tồn tại $w$ tách được cách tô này, thì với mọi $i\ge 1$ ta có
\[
  w^Ta_ix_i\ge 0.
\]
Suy ra
\[
  w^Tx_0=\sum_{i=1}^{n+1}w^Ta_ix_i\ge 0.
\]
Nhưng $y_0=-1$ lại đòi hỏi $f_w(x_0)=w^Tx_0<0$, mâu thuẫn. Vậy với
$n+2$ điểm luôn có một cách tô không thể tách bởi siêu phẳng $n-1$ chiều.

\section{1004 Another String}

Định nghĩa $F[i,j]$ là độ dài lớn nhất của hai chuỗi con bắt đầu lần lượt tại
$i,j$ trong $S$ sao cho chúng thỏa $k$-matching. Nói cách khác, $F[i,j]$ là
giá trị lớn nhất của $L$ để
\[
  S[i,i+L-1]\quad\text{và}\quad S[j,j+L-1]
\]
thỏa $k$-matching.

Để tính mọi $F[i,j]$ với $1\le i<j\le n$, ta liệt kê hiệu hai đầu trái
$d\in[1,n-1]$. Với mỗi $d$, nhờ tính đơn điệu của $k$-matching, dùng hai con
trỏ để tính toàn bộ $F[i,i+d]$ trong thời gian tuyến tính. Vì vậy mọi
$F[i,j]$ được tính trong $O(n^2)$.

Tiếp theo xét mọi cách chia chuỗi. Gọi phép chia tại $t$ là
\[
  A=S[1,t-1],\qquad B=S[t,n].
\]
Với phép chia này, đáp án cần thống kê là
\[
  \sum_{i=1}^{t-1}\sum_{j=t}^{n}G[i,j],
  \qquad
  G[i,j]=\min(F[i,j],t-i).
\]
Ta liệt kê $t$ từ lớn xuống nhỏ và duy trì đóng góp của $G$. Vì
$G[i,j]\le t-i$, với mỗi $i$ duy trì mảng $\operatorname{count}_i$ để đếm số
giá trị $G[i,j]$ đang đóng góp; $\operatorname{count}_i[x]$ là số giá trị có
đóng góp bằng $x$. Đồng thời duy trì $\operatorname{max}_i$, giá trị đóng góp
lớn nhất hiện tại của hàng $i$.

Mỗi khi $t$ giảm, ta cập nhật các mảng đếm và các giá trị lớn nhất để tính
thay đổi của đáp án. Trong quá trình này $\operatorname{max}_i$ chỉ giảm; số
giá trị $G[i,j]$ mới sinh đóng góp mỗi lần là $O(n)$ và giá trị của chúng
cũng không vượt quá $t-i$. Tổng thời gian là $O(n^2)$.

\section{1005 Random Walk on Graph}

Đặt
\[
  \lambda_i=W_{i,i},\qquad d_i=\sum_{j=1}^{n} W_{i,j}.
\]
Ta có phương trình chuyển:
\[
  A_{i,i}
  =
  \frac{\lambda_i}{\lambda_i+d_i}
  +
  \sum_{\substack{j=1\\j\ne i}}^{n}
  \frac{W_{i,j}}{\lambda_i+d_i}A_{j,i},
\]
và với $i\ne j$,
\[
  A_{i,j}
  =
  \sum_{\substack{k=1\\k\ne i}}^{n}
  \frac{W_{i,k}}{\lambda_i+d_i}A_{k,j}.
\]
Gọi $W$ là ma trận không có khuyên, $\Lambda$ là ma trận đường chéo của các
khuyên, và $D$ là ma trận bậc của $W$. Khi đó dạng ma trận là
\[
  \bigl(I-(\Lambda+D)^{-1}W\bigr)A=(\Lambda+D)^{-1}\Lambda.
\]
Các bước là: lập phương trình chuyển, viết thành dạng ma trận, rồi giải hệ
ma trận bằng mẫu có sẵn. Sau khi tính được ma trận, chỉ cần dùng phép nghịch
đảo ma trận.

\section{1006 Cute Tree}

Chỉ cần mô phỏng đúng theo đề bài và in số nút của cây.

\section{1007 Banzhuan}

Với chi phí lớn nhất, rõ ràng cần lấp đầy toàn bộ không gian
$n\times n\times n$. Để chi phí lớn nhất, mỗi viên gạch đều được thả từ vị
trí cao nhất. Do đó
\[
  \sum_{x=1}^{n}\sum_{y=1}^{n}\sum_{z=1}^{n} x y^2 n
  =
  n^2\cdot\frac{(n+1)n}{2}\cdot\frac{n(n+1)(2n+1)}{6}.
\]

Với chi phí nhỏ nhất, do trọng lực, mặt đáy phải được phủ kín. Để thỏa hai
hình chiếu còn lại, cần thêm một bức tường đứng tăng theo $x$ và giảm theo
$y$. Vì vậy chi phí nhỏ nhất là
\[
\begin{aligned}
  &\sum_{x=1}^{n}\sum_{y=1}^{n}xy
  +\sum_{y=1}^{n}\sum_{z=2}^{n}(n+1-y)y^2z\\
  &\quad =
  \frac{(n+1)^2n^2}{4}
  +\frac{(n+2)(n-1)}{2}
  \left(
    \frac{n(n+1)^2(2n+1)}{6}
    -\frac{n^2(n+1)^2}{4}
  \right).
\end{aligned}
\]

\section{1008 Set Counting}

Chú ý rằng đề bài yêu cầu một xác suất có điều kiện.

Dùng tiền tố tổng trên siêu lập phương để tính, với mỗi tập $T$, có bao nhiêu
siêu tập $U$ của nó. Gọi giá trị đó là $\operatorname{cnt}[T]$. Khi đó đáp
án ban đầu là
\[
  \sum_{S\in[0,2^n)}\sum_{T\in[0,2^n)}
  \frac{\operatorname{cnt}[T\vee S]}{\operatorname{cnt}[S]}.
\]
Biểu thức $T\vee S$ lại có thể được tối ưu bằng một lần tiền tố tổng cao
chiều nữa. Gọi kết quả là $\operatorname{cnt0}[T]$. Khi đó
\[
  \sum_{S\in[0,2^n)}
  \frac{\operatorname{cnt0}[S]\cdot 2^{\operatorname{bit}(S)}}
       {\operatorname{cnt}[S]}.
\]
Ở đây $\operatorname{bit}(S)$ là số bit $1$ trong $S$. Độ phức tạp là
$O(n2^n)$.

\section{1009 Array Counting}

Cách thô có độ phức tạp $O(kn\log n)$. Cách dùng cây đoạn theo trọng số hoặc
cây Fenwick được lược bỏ trong lời giải gốc.

\paragraph{Cách 3: chia theo số lần xuất hiện.}
Khi $k$ lớn, chia các giá trị theo số lần xuất hiện: ít nhất $\sqrt n$ và
nhỏ hơn $\sqrt n$. Đếm số lần xuất hiện của mỗi giá trị bằng bucket, bảng
băm hoặc tập, rồi sắp theo thứ tự giảm dần.

Nếu một giá trị xuất hiện ít nhất $\sqrt n$ lần, số loại như vậy nhiều nhất
$\sqrt n$. Phần này dùng cách cây đoạn hoặc Fenwick, có độ phức tạp
$\sqrt n\cdot n\log n$.

Nếu một giá trị xuất hiện ít hơn $\sqrt n$ lần, xét các vị trí xuất hiện của
nó. Giả sử giá trị đó xuất hiện $t_i$ lần. Liệt kê các lựa chọn
$l_i,r_i,t_i$ rồi tính những vị trí hợp lệ ở trái và phải; một giá trị xử lý
được trong $O(t_i^2)$. Tổng
\[
  \sum_i t_i^2
  \le
  \sqrt n\sum_i t_i
  =
  n\sqrt n.
\]

\paragraph{Cách 4: cân bằng hai thuật toán.}
Đặt ngưỡng chia khối là $d$. Với các giá trị xuất hiện ít nhất $d$ lần, số
loại nhiều nhất $n/d$, độ phức tạp là
\[
  \frac{n}{d}\cdot n\log n.
\]
Với các giá trị xuất hiện ít hơn $d$ lần, tổng chi phí là $nd$. Chọn
\[
  d=\sqrt{n\log n}
\]
ta được độ phức tạp
\[
  O(n^{3/2}\sqrt{\log_2 n}).
\]

\paragraph{Cách 5: cây Fenwick.}
Độ phức tạp là $O(n\log_2 n)$. Đếm số lần xuất hiện của mỗi giá trị, rồi
lần lượt xét từng giá trị và hợp nhất ảnh hưởng của hai khối kề nhau cùng
giá trị. Đáp án có thể nhìn như: biến mỗi vị trí thành $1$ hoặc $-1$, lấy
tiền tố, rồi đếm trước mỗi tiền tố có bao nhiêu tiền tố nhỏ hơn nó.

Ví dụ mảng ban đầu
\[
  1,1,2,1
\]
được đổi thành
\[
  1,1,-1,1,
\]
tiền tố là
\[
  0,1,2,1,2,
\]
và các đóng góp lần lượt là $1,2,1,3$.

Nếu tiền tố hiện tại lớn hơn một tiền tố trước đó, thì đoạn giữa hai tiền tố
có tổng dương:
\[
  \operatorname{sum}[\operatorname{now}]
  -\operatorname{sum}[\operatorname{last}]>0.
\]
Vì vậy đoạn $\operatorname{last}+1\ldots\operatorname{now}$ hợp lệ. Phần còn
lại là tính đóng góp của từng vị trí $1$ đối với các vị trí $-1$ trước nó.
Nếu có $k$ loại, loại thứ $i$ xuất hiện $t_i$ lần, mỗi lần tốn $\log n$, thì
tổng độ phức tạp là
\[
  \sum_i t_i\log n=O(n\log n).
\]

\paragraph{Cách 6.}
Có thể tối ưu cách 5 xuống $O(n)$. Gọi $\operatorname{sum}$ là tiền tố hiện
tại, $f_1[\operatorname{sum}]$ là số điểm có tiền tố bằng giá trị này. Với
mỗi điểm, đóng góp hiện tại là tổng các $f_1$ ở tiền tố nhỏ hơn
$\operatorname{sum}$. Nếu tiếp theo có một đoạn liên tiếp toàn $-1$, ta nhảy
ngay đến cuối đoạn đó. Dùng thêm một mảng hiệu $f_2$: đánh dấu $-1$ ở đầu
đoạn và $+1$ ở cuối đoạn. Khi đi đến $i$, nếu $f_2[\operatorname{sum}]$ khác
$0$, dùng nó cập nhật $f_1$, đồng thời cập nhật $f_2[\operatorname{sum}+1]$.
Chi tiết cài đặt được lược bỏ.

\section{1010 Gumbel Softmax}

Với $i=1,\ldots,k$,
\[
  y_i=
  \frac{\exp\left((x_i+g_i-(x_k+g_k))/\tau\right)}
       {\sum_{j=1}^{k}\exp\left((x_j+g_j-(x_k+g_k))/\tau\right)}.
\]
Đặt
\[
  u_i=x_i+g_i-(x_k+g_k),\qquad i=1,\ldots,k-1.
\]
Ở đây $g_i$ là biến ngẫu nhiên trong mã giả, tuân theo phân bố Gumbel. Mật
độ của phân bố Gumbel là
\[
  f(z,\mu)=e^{\mu-z-e^{\mu-z}}.
\]

Trước hết tính mật độ đồng thời của $u$:
\[
\begin{aligned}
  p(u_1,\ldots,u_{k-1})
  &=
  \int_{-\infty}^{\infty}
  dg_k\,p(g_k)\prod_{i=1}^{k-1}p(u_i\mid g_k)\\
  &=
  \int_{-\infty}^{\infty}
  dg_k\,f(g_k,0)\prod_{i=1}^{k-1}f(x_k+g_k,x_i-u_i).
\end{aligned}
\]
Sau khi đổi biến và rút gọn, nhận được
\[
  p(u_1,\ldots,u_{k-1})
  =
  \Gamma(k)
  \left(\prod_{i=1}^{k}e^{x_i-u_i}\right)
  \left(\sum_{i=1}^{k}e^{x_i-u_i}\right)^{-k},
\]
trong đó quy ước $u_k=0$.

Tiếp theo cần mật độ đồng thời của $y$. Ta có
\[
  y_k=1-\sum_{j=1}^{k-1}y_j,
\]
và
\[
  h_i(u)=
  \frac{e^{u_i/\tau}}{1+\sum_{j=1}^{k-1}e^{u_j/\tau}}.
\]
Dùng đổi biến:
\[
  p(y_{1:k})
  =
  p(h^{-1}(y_{1:k-1}))
  \det\left(
    \frac{\partial h^{-1}(y_{1:k-1})}{\partial y_{1:k-1}}
  \right).
\]
Nghịch đảo là
\[
  h^{-1}_i(y)=\tau(\ln y_i-\ln y_k).
\]
Jacobian có dạng
\[
  \tau\left(\operatorname{diag}\left(\frac1{y_{1:k-1}}\right)
  +\frac1{y_k}\mathbf{1}\mathbf{1}^T\right),
\]
và định thức là
\[
  \tau^{k-1}\prod_{j=1}^{k} y_j^{-1}.
\]
Do đó
\[
  p(y_1,\ldots,y_k)
  =
  \Gamma(k)\tau^{k-1}
  \left(\sum_{i=1}^{k}\frac{e^{x_i}}{y_i^\tau}\right)^{-k}
  \prod_{i=1}^{k}\frac{e^{x_i}}{y_i^{\tau+1}}.
\]

\section{1011 Stone Game}

\paragraph{Trường hợp $x=y$.}
Dễ chứng minh
\[
  \operatorname{sg}_i=\left\lfloor\frac{i}{2x}\right\rfloor x+i\bmod x.
\]
Khi đó bài toán trở thành Nim: nếu
\[
  \bigoplus_{i=1}^{n}\operatorname{sg}_{a_i}>0
\]
thì người đi trước thắng, ngược lại thua.

Sau đây đặt $\operatorname{sg}_i=i$, tức giá trị SG của Nim.

\paragraph{Trường hợp $x<y$.}
Giả sử
\[
  a_1\le a_2\le\cdots\le a_n.
\]
Nếu $a_n<x$, bài toán là Nim, nên chỉ cần xét $a_n\ge x$.

Nếu người đi trước có thể qua một thao tác làm cho
\[
  \max_i a_i<x,
\]
thì thắng thua do $\bigoplus_i a_i$ quyết định. Nếu tìm được số dương $z$
sao cho
\[
  \bigoplus_{i=1}^{n}a_i\oplus a_n\oplus(a_n-z)=0,
\]
thì người đi trước thắng; nếu không, người đi trước phải xét chiến lược khác.

Khi người đi trước không thể đưa mọi đống xuống dưới $x$ trong một thao tác,
phân tích trạng thái sau lượt đầu theo ba trường hợp.

\paragraph{1. $a_{n-1}<x$ và $a_n\ge x$.}
Nếu xor hiện tại khác $0$, tồn tại một đống $j$ và số dương $z$ sao cho
\[
  \bigoplus_i a_i\oplus a_j\oplus(a_j-z)=0.
\]
Nếu $z\ne y$, người đi sau trực tiếp lấy $z$ viên để đưa người đi trước vào
trạng thái xor bằng $0$. Nếu $z=y$, khi đó $j=n$; người đi sau lấy
$y-x$ viên từ $a_n$, khiến người đi trước không thể lấy $x$ viên từ đống thứ
$n$.

Nếu xor hiện tại bằng $0$, giả sử người đi sau lấy $x$ viên từ $a_n$. Theo
tính chất của Nim và $a_{n-1}<x<y$, người đi trước chắc chắn có nước đáp trả.
Người đi sau chọn đúng nước đáp trả đó để khiến người đi trước rơi vào tình
huống không thể lấy $x$ viên từ đống thứ $n$.

Vì vậy mọi cục diện người đi trước có thể tạo ra đều là
\[
  \bigoplus_i a_i=0
  \quad\text{hoặc}\quad
  \bigoplus_i a_i\oplus a_n\oplus(a_n-x)=0,
\]
và người đi sau thắng.

\paragraph{2. $a_{n-2}<x$, $a_{n-1}\ge x$, $a_n\ge x$.}
Người đi trước sẽ không tự lấy để làm $a_{n-1}<x$ hoặc $a_n<x$, và người đi
sau cũng vậy. Cuối cùng còn hai đống $a_{n-1}=x$ và $a_n=x$. Nếu đến lượt
người đi trước, người đi sau đối xứng; nếu đến lượt người đi sau, người đi
sau lấy hết một đống $x$, còn lại một đống $x$ mà người đi trước không thể
lấy hết trong một lượt. Người đi sau thắng.

\paragraph{3. $a_{n-2}\ge x$, $a_{n-1}\ge x$, $a_n\ge x$.}
Người đi sau chỉ cần sau mỗi lượt vẫn duy trì trường hợp này, buộc người đi
trước tự đi vào trường hợp 2. Người đi trước sẽ không chủ động vào trường hợp
2; cuối cùng còn ba đống $x$, người đi sau lấy hết một đống rồi đối xứng.

Tóm lại, khi $x<y$, người đi trước thắng khi có thể qua một thao tác làm
$\max_i a_i<x$ và thỏa điều kiện xor ở trên; nếu không thì thua.

\paragraph{Trường hợp $x>y$.}
Nếu $a_n<y$, bài toán là Nim. Nếu $a_n\ge y$, tình huống đối ứng với ba
trường hợp đã xét khi $x<y$, và người đi trước thắng.

\section{1012 Matrix Counting}

Chú ý
\[
\begin{aligned}
  \operatorname{ans}[i]
  &=
  \left[
    \sum_{i=1}^{n}\sum_{j=1}^{r}S_{ij}=i
  \right]\\
  &=
  \left[
    \sum_{i=1}^{n}\sum_{j=1}^{n}\sum_{k=1}^{r} A_{ik}B_{kj}=i
  \right]\\
  &=
  \left[
    \sum_{k=1}^{r}
    \left(\sum_{i=1}^{n} A_{ik}\right)
    \left(\sum_{j=1}^{n} B_{kj}\right)=i
  \right].
\end{aligned}
\]
Gọi $F[q]$ là số phương án để
\[
  \sum_{i=1}^{n}A_{ik}=q,
\]
và $G[q]$ tương tự cho $B$. Rõ ràng $F[q]=G[q]$.

Để tính $F[x]$, xét hàm sinh
\[
  g(x)=(1+x+x^2+\cdots+x^m)^n,
\]
khi đó
\[
  F[q]=[x^q]g(x).
\]
Có hai cách tính. Với $q=0,1,\ldots,m$, hệ số bậc $q$ tương đương số nghiệm
nguyên không âm của
\[
  x_1+x_2+\cdots+x_n=q,
\]
nên bằng
\[
  [x^q]g(x)=\binom{n+q-1}{n-1}.
\]
Ngoài ra có thể dùng lũy thừa đa thức nhanh bằng cách lấy $\ln$ rồi $\exp$.

Gọi $H[i]$ là số phương án để
\[
  \left(\sum_i A_{ik}\right)\left(\sum_j B_{kj}\right)=i.
\]
Ta có thể liệt kê $x,y$ và cộng
\[
  H[xy]\mathrel{+}=F[x]G[y]
\]
để thu được các giá trị $H[i]$ với $i=0,1,\ldots,m$.

Lưu ý rằng các giá trị như
\[
  H[0]\mathrel{+}=F[0]G[y],\qquad y=m+1,\ldots,mn
\]
chưa được tính trong phép liệt kê trên. Do
\[
  \sum_{i=0}^{mn}F[i]=\sum_{i=0}^{mn}G[i]=(m+1)^n,\qquad F[0]=G[0]=1,
\]
cần hiệu chỉnh
\[
  H[0]=(m+1)^n-1.
\]
Đặt hàm sinh
\[
  h(x)=(H[0]+H[1]x+H[2]x^2+\cdots+H[m]x^m)^r.
\]
Khi đó
\[
  \operatorname{ans}[i]=[x^i]h(x).
\]
Ta dùng lũy thừa đa thức nhanh để tính $h(x)$. Vì $r$ rất lớn, dùng định lý
Euler để xử lý số mũ. Độ phức tạp là $O(m\log m)$.

\section{1013 Tree Diameter}

Bài toán tương đương với: sau khi thao tác trên $n$ điểm, làm cho đường kính
của cây ngắn nhất có thể.

Nhị phân đáp án. Mỗi lần kiểm tra, dùng giá trị giữa làm ngưỡng để hạn chế
chuyển trạng thái trong quy hoạch động trên cây.

Với cây con gốc $u$, duy trì hai trạng thái. $dp[u][0]$ biểu diễn trường hợp
giá trị power của $u$ được cho cho con; khi đó nó lưu độ dài chuỗi dài nhất
từ một lá trong cây con của $u$ đến $u$. $dp[u][1]$ biểu diễn trường hợp
power của $u$ được cho cho cha, cũng lưu độ dài chuỗi dài nhất từ lá đến
$u$.

Với chuyển trạng thái của $dp[u][1]$, chỉ cần bảo đảm mọi cặp con đi qua
$u$ đều có tổng độ dài không vượt ngưỡng, rồi ghi lại chuỗi dài nhất. Với
$dp[u][0]$, liệt kê việc power của $u$ đóng góp cho chuỗi dài nhất hoặc chuỗi
dài thứ hai; trong số các lựa chọn thỏa mọi tổng độ dài qua hai con không
vượt ngưỡng, lấy chuỗi dài nhất nhỏ nhất. Cuối cùng chỉ cần kiểm tra trạng
thái tại gốc có hợp lệ hay không.

Độ phức tạp là
\[
  O(n\log(n\cdot \max w)).
\]

\end{document}
